\section{从语言模型到助手}

\subsection{预训练语言模型的局限}

预训练语言模型学会了\textbf{补全文本}，但这不等于\textbf{帮助用户}。语言建模目标与"满足人类偏好"之间存在根本性的不对齐。

\begin{knowledgebox}{预训练学到了什么}
预训练通过大量文本学到了丰富的知识：事实、语法、共指消解、情感分析、基本推理、数学、编程等。但它学到的是"下一个最可能的 token"，而不是"用户最想要的回答"。
\end{knowledgebox}

\subsection{人类偏好的多样性}

真实用户偏好并不是单一的奖赏函数，而是由多个维度组成：任务完成度、创造性、礼貌、可解释性、甚至审美。Archit 在讲座中称\textit{``偏好是一个多模态的向量空间，而非单一标量''}，并强调在构建信号时一定要分层采集。

我们可以按来源将偏好分为：
\begin{itemize}
    \item \textbf{人类比较}：用户或标注者给出回答队列，挑选最优者。
    \item \textbf{可验证指标}：数学题目、编程任务、事实一致性等可以自动打分。
    \item \textbf{行为数据}：点击率、重写率、跳出率等隐含用户满意度。
    \item \textbf{治理约束}：安全、合规要求需要额外的 guardrail 信号。
\end{itemize}

\begin{knowledgebox}{偏好矩阵建模的洞察}
把多个偏好信号组合成 \textit{hierarchical preference matrix}，能够明确每种信号的作用范围，并在优化时以加权方式对齐不同目标。
\end{knowledgebox}

\begin{importantbox}{多信号偏好可解释化}
把 signal 分成\textit{primary preference}（核心满意度）和\textit{secondary constraint}（安全、bias）两类，分别用 reward model 和 guardrail filter 处理，能减少 reward hacking 并提升治理透明度。
\end{importantbox}

\subsection{本章小结}

预训练 LM 具有丰富的知识，但要变成助手必须从指令微调、偏好信号拆分、治理信号嵌入等多个维度对齐，才能避免单一 reward 走偏。

